\documentclass[12pt,a4paper]{article}

\usepackage[utf8]{inputenc}
\usepackage[vietnamese]{babel}
\usepackage{amsmath,amssymb,amsthm}
\usepackage{geometry}
\usepackage{enumitem}
\usepackage[most]{tcolorbox}
\usepackage{hyperref}
\usepackage{fancyhdr}
\usepackage{tikz}
\usepackage{eso-pic}
\usepackage{xcolor}

\geometry{a4paper, margin=2cm, bottom=2.5cm}
\hypersetup{colorlinks=true, linkcolor=blue!70!black}
\setlist[itemize]{topsep=4pt, itemsep=2pt}
\setlist[enumerate]{topsep=4pt, itemsep=2pt}

\AddToShipoutPictureFG{%
  \AtPageCenter{%
    \begin{tikzpicture}[remember picture, overlay]
      \node[rotate=45, scale=1, opacity=0.06]{\fontsize{60}{72}\selectfont\bfseries\sffamily Đặng Tấn Quí};
    \end{tikzpicture}%
  }%
}

\pagestyle{fancy}
\fancyhf{}
\fancyhead[L]{\small\textit{Hình học giải tích}}
\fancyhead[R]{\small\textit{Đặng Tấn Quí}}
\fancyfoot[C]{\thepage}
\fancyfoot[L]{\footnotesize\textcopyright\ Đặng Tấn Quí}
\fancyfoot[R]{\footnotesize SĐT: 0377\,122\,210}
\renewcommand{\headrulewidth}{0.4pt}
\renewcommand{\footrulewidth}{0.4pt}

\fancypagestyle{plain}{\fancyhf{}\fancyfoot[C]{\thepage}\fancyfoot[L]{\footnotesize\textcopyright\ Đặng Tấn Quí}\fancyfoot[R]{\footnotesize SĐT: 0377\,122\,210}\renewcommand{\headrulewidth}{0pt}\renewcommand{\footrulewidth}{0.4pt}}

\newtcolorbox{dinhnghia}[1][]{enhanced, breakable, colback=blue!3!white, colframe=blue!60!black, fonttitle=\bfseries, title={Định nghĩa}, #1}
\newtcolorbox{dinhly}[1][]{enhanced, breakable, colback=green!3!white, colframe=green!50!black, fonttitle=\bfseries, title={Định lý}, #1}
\newtcolorbox{tinhchat}[1][]{enhanced, breakable, colback=yellow!5!white, colframe=orange!50!black, fonttitle=\bfseries, title={Tính chất}, #1}
\newtcolorbox{phuongphap}[1][]{enhanced, breakable, colback=gray!5!white, colframe=gray!50!black, fonttitle=\bfseries, title={Phương pháp}, #1}
\newtcolorbox{luuy}[1][]{enhanced, breakable, colback=red!3!white, colframe=red!50!black, fonttitle=\bfseries, title={Lưu ý}, #1}
\newtcolorbox{congthuc}[1][]{enhanced, breakable, colback=cyan!3!white, colframe=cyan!50!black, fonttitle=\bfseries, title={Công thức}, #1}
\newtcolorbox{vidu}[1][]{enhanced, breakable, colback=violet!3!white, colframe=violet!50!black, fonttitle=\bfseries, title={Ví dụ}, #1}

\begin{document}

\thispagestyle{empty}
\begin{tikzpicture}[remember picture, overlay]
  \fill[blue!80!black] (current page.south west) rectangle (current page.north east);
  \node[anchor=west, white, font=\fontsize{26}{32}\bfseries\selectfont]
    at ([xshift=3cm, yshift=-7.5cm]current page.north west) {HÌNH HỌC GIẢI TÍCH};
  \node[anchor=west, white, font=\fontsize{18}{24}\selectfont]
    at ([xshift=3cm, yshift=-9cm]current page.north west) {Đường thẳng, mặt phẳng, mặt bậc hai --- Năm 1, HK2};
  \node[anchor=west, white, font=\Large\bfseries] at ([xshift=3cm, yshift=-13cm]current page.north west) {Biên soạn:};
  \node[anchor=west, yellow!80!white, font=\fontsize{22}{28}\bfseries\selectfont] at ([xshift=3cm, yshift=-14.5cm]current page.north west) {ĐẶNG TẤN QUÍ};
  \node[anchor=south east, white, font=\Large\bfseries] at ([xshift=-2cm, yshift=3cm]current page.south east) {\the\year};
\end{tikzpicture}

\newpage
\tableofcontents
\newpage

\section{Vectơ trong không gian}

\subsection{Vectơ và các phép toán}

\begin{dinhnghia}
  \textbf{Vectơ} $\vec{u} = (u_1, u_2, u_3)$. Độ dài: $|\vec{u}| = \sqrt{u_1^2 + u_2^2 + u_3^2}$.
\end{dinhnghia}

\begin{congthuc}[title={Tích vô hướng}]
  $\vec{u} \cdot \vec{v} = u_1 v_1 + u_2 v_2 + u_3 v_3 = |\vec{u}||\vec{v}|\cos\varphi$.
\end{congthuc}

\begin{congthuc}[title={Tích có hướng}]
  $\vec{u} \times \vec{v} = \begin{vmatrix} \vec{i} & \vec{j} & \vec{k} \\ u_1 & u_2 & u_3 \\ v_1 & v_2 & v_3 \end{vmatrix}$. $|\vec{u} \times \vec{v}| = |\vec{u}||\vec{v}||\sin\varphi|$.
\end{congthuc}

\begin{congthuc}[title={Tích hỗn tạp}]
  $[\vec{u}, \vec{v}, \vec{w}] = \vec{u} \cdot (\vec{v} \times \vec{w}) = \det(\vec{u}, \vec{v}, \vec{w})$. Thể tích hình hộp $= |[\vec{u}, \vec{v}, \vec{w}]|$.
\end{congthuc}

\section{Đường thẳng trong không gian}

\begin{dinhnghia}
  Đường thẳng qua $M_0(x_0, y_0, z_0)$ có vectơ chỉ phương $\vec{u} = (a, b, c)$:
  \begin{itemize}
    \item \textbf{Tham số}: $x = x_0 + at$, $y = y_0 + bt$, $z = z_0 + ct$.
    \item \textbf{Chính tắc}: $\dfrac{x - x_0}{a} = \dfrac{y - y_0}{b} = \dfrac{z - z_0}{c}$ (khi $abc \ne 0$).
  \end{itemize}
\end{dinhnghia}

\begin{phuongphap}[title={Vị trí tương đối hai đường thẳng}]
  Xét $\vec{u}_1$, $\vec{u}_2$ (VTCP) và $\overrightarrow{M_1M_2}$. Dùng tích có hướng và tích hỗn tạp để xét song song, cắt nhau, chéo nhau.
\end{phuongphap}

\begin{congthuc}[title={Khoảng cách từ điểm đến đường thẳng}]
  $d(M, \Delta) = \dfrac{\bigl|\overrightarrow{M_0M} \times \vec{u}\bigr|}{|\vec{u}|}$, với $M_0 \in \Delta$, $\vec{u}$ là VTCP.
\end{congthuc}

\section{Mặt phẳng}

\begin{dinhnghia}
  Mặt phẳng $(P)$ qua $M_0(x_0, y_0, z_0)$ có VTPT $\vec{n} = (A, B, C)$:
  \[
    A(x - x_0) + B(y - y_0) + C(z - z_0) = 0 \quad \Leftrightarrow \quad Ax + By + Cz + D = 0.
  \]
\end{dinhnghia}

\begin{congthuc}[title={Khoảng cách từ điểm đến mặt phẳng}]
  $d(M, (P)) = \dfrac{|Ax_0 + By_0 + Cz_0 + D|}{\sqrt{A^2 + B^2 + C^2}}$.
\end{congthuc}

\begin{phuongphap}[title={Góc giữa hai mặt phẳng}]
  Góc giữa $(P_1)$ và $(P_2)$ bằng góc giữa $\vec{n}_1$ và $\vec{n}_2$: $\cos\varphi = \dfrac{|\vec{n}_1 \cdot \vec{n}_2|}{|\vec{n}_1||\vec{n}_2|}$.
\end{phuongphap}

\section{Mặt bậc hai}

\begin{dinhnghia}
  \textbf{Mặt cầu} tâm $I(a, b, c)$, bán kính $R$: $(x-a)^2 + (y-b)^2 + (z-c)^2 = R^2$.
\end{dinhnghia}

\begin{dinhnghia}
  \textbf{Ellipsoid}: $\dfrac{x^2}{a^2} + \dfrac{y^2}{b^2} + \dfrac{z^2}{c^2} = 1$.
\end{dinhnghia}

\begin{dinhnghia}
  \textbf{Mặt trụ}: phương trình thiếu một biến, ví dụ $x^2 + y^2 = R^2$ (trụ tròn).
\end{dinhnghia}

\begin{dinhnghia}
  \textbf{Mặt nón}: $z^2 = \dfrac{x^2}{a^2} + \dfrac{y^2}{b^2}$.
\end{dinhnghia}

\begin{dinhnghia}
  \textbf{Paraboloid elliptic}: $z = \dfrac{x^2}{a^2} + \dfrac{y^2}{b^2}$. \textbf{Paraboloid hyperbolic} (yên ngựa): $z = \dfrac{x^2}{a^2} - \dfrac{y^2}{b^2}$.
\end{dinhnghia}

\end{document}
